\chapter{Antiderivatives}\label{ch:anti}

Integration is taught as a bag of tricks, and the relearning is to see that the bag has a
structure: every technique is a derivative rule read backwards, and every answer can be checked
by reading it forwards again. The engine takes that literally. Its integrator is a guesser; its
theorem is about the checker.

\section{Antiderivatives and the fundamental theorem}

\begin{definition}
  $F$ is an antiderivative of $f$ on an interval $I$ if $F'(x) = f(x)$ for all $x \in I$.
\end{definition}

\begin{theorem}
  If $F$ and $G$ are antiderivatives of $f$ on an interval $I$, then $F - G$ is constant on $I$.
\end{theorem}
\begin{proof}
  $(F - G)' = 0$ on $I$. By the mean value theorem, for $a < b$ in $I$ there is $c \in (a, b)$ with
  $(F-G)(b) - (F-G)(a) = (F-G)'(c)(b - a) = 0$.
\end{proof}

The mean value theorem is the step, and it is the reason $I$ must be an interval: on
$(-\infty, 0) \cup (0, \infty)$, the function that is $1$ on the left and $2$ on the right has
zero derivative and is not constant. The constant of integration is this theorem, and the
engine omits it deliberately: it returns \emph{one} antiderivative, and its claim is about
that one.

\begin{theorem}[Fundamental theorem of calculus]
  If $f$ is continuous on $[a, b]$ then $F(x) = \int_a^x f$ is an antiderivative of $f$, and for
  any antiderivative $G$, $\int_a^b f = G(b) - G(a)$.
\end{theorem}

The second part is the first part with the previous theorem. The book does not need the
integral itself, only the equation $F' = f$; that equation is what ``antiderivative'' means and
what the engine checks.

\section{The table, backwards}

Every derivative rule of Chapters~\ref{ch:dderiv} and~\ref{ch:power} read right to left is an
integration rule.

\begin{center}
\begin{tabular}{ll}
  \toprule
  Derivative & Antiderivative \\
  \midrule
  $(x^{n+1}/(n+1))' = x^n$ & $\int x^n\,dx = x^{n+1}/(n+1)$, $n \neq -1$ \\
  $(\ln x)' = 1/x$ & $\int x^{-1}\,dx = \ln x$, $x > 0$ \\
  $(e^x)' = e^x$ & $\int e^x\,dx = e^x$ \\
  $(-\cos x)' = \sin x$ & $\int \sin x\,dx = -\cos x$ \\
  $(\sin x)' = \cos x$ & $\int \cos x\,dx = \sin x$ \\
  $(x \ln x - x)' = \ln x$ & $\int \ln x\,dx = x\ln x - x$ \\
  $(-\ln \cos x)' = \tan x$ & $\int \tan x\,dx = -\ln \cos x$ \\
  \bottomrule
\end{tabular}
\end{center}

Linearity of the derivative gives linearity of the antiderivative: $\int (f + g) = \int f +
\int g$ and $\int cf = c \int f$.

\section{Substitution}

\begin{theorem}[Substitution]
  If $G' = g$ and $u$ is differentiable, then
  \[ \int g(u(x))\,u'(x)\,dx = G(u(x)). \]
\end{theorem}
\begin{proof}
  The chain rule: $(G \circ u)' = (G' \circ u)\, u' = (g \circ u)\, u'$.
\end{proof}

The special case $u = ax + b$ gives $\int g(ax + b)\,dx = G(ax+b)/a$, since $u' = a$ is a
constant that can be divided out. The general case is the one to learn to \emph{see}: in
$\int 2x \cos(x^2)\,dx$ the factor $2x$ is the derivative of the inner $x^2$, so the answer is
$\sin(x^2)$. When the derivative is present only up to a constant, as in $\int x \cos(x^2)\,dx$,
the constant comes out: $\frac{1}{2}\sin(x^2)$.

\section{The finder}

The engine's antidifferentiator is a direct transcription of the last two sections, and nothing
about it is proved.

\begin{minted}{lean}
def table (g : String) (u : Expr) : Option (Expr × String) :=
  match g with
  | "sin" => some (neg (.fn "cos" [u]), "$\\int \\sin u \\, du = -\\cos u$")
  | "cos" => some (.fn "sin" [u], "$\\int \\cos u \\, du = \\sin u$")
  | "exp" => some (.fn "exp" [u], "$\\int e^u \\, du = e^u$")
  | "ln" => some (sub (.mul [u, .fn "ln" [u]]) u, "$\\int \\ln u \\, du = u \\ln u - u$")
  | "tan" => some (neg (.fn "ln" [.fn "cos" [u]]), "$\\int \\tan u \\, du = -\\ln \\cos u$")
  | _ => none

def powerRule (u n : Expr) : Expr × String :=
  if n.isNumEq Q.minusOne then (.fn "ln" [u], "$\\int u^{-1} \\, du = \\ln u$")
  else match n with
    | .num q => let n1 := q.add Q.one; (.mul [.num (Q.one.div n1), .pow u (.num n1)], "…")
    | _ => let n1 := .add [n, Expr.one]; (.mul [.pow n1 Expr.minusOne, .pow u n1], "… (assuming $n \\neq -1$)")

def substitute (x : String) (f G : Expr) (a : Expr) : Expr × Array Step :=
  if a.isOne then (G, #[]) else (Expr.div G a, #[⟨"int.linear-substitution", "…", [], integral f x, Expr.div G a, none⟩])
\end{minted}

The main function \lc{anti} dispatches on the shape of the integrand: a term free of $x$ is a
constant; a sum integrates termwise; constant factors come out; a power $u^n$ or a table
function of $u$ with $u$ linear in $x$ uses the table and \lc{substitute}; and a product tries
the general substitution --- one factor is $H'(u)$ for a table or power $H$, and the remaining
factors, divided by $u'$, must be free of $x$:

\begin{minted}{lean}
let trySubst (i : Nat) (w G : Expr) (why : String) : Option (Expr × Array Step) := do
  if !w.dependsOn x then none else
  let w' ← simp (D w x)
  let ratio ← simp (Expr.div (without es i) w')
  if ratio.dependsOn x then none else
  let F := .mul [ratio, G]
  return (F, #[step "int.substitution" "…" F])
\end{minted}

The derivative $w'$ is computed by the caller's normalizer --- the pipeline of
Chapter~\ref{ch:dderiv} --- so the finder never differentiates on its own. Its steps are
recorded with the status \status{checked}: a guess whose result a later step verifies.

\section{The checker}

\begin{minted}{lean}
def cmdIntegrate (norm : Norm) : PlainRule :=
  { name := "cmd.integrate", apply := fun e => Option.map checked <|
      match e with
      | .fn "integrate" [f, .var x] =>
        match Anti.anti (fun e => (norm e).toOption.map (·.1)) x 3 f with
        | none => some (refuse "integrate: no antiderivative … found")
        | some (F₀, steps) =>
          match norm F₀ with
          | .ok (F, _) =>
            match norm (D F x) with
            | .ok (g, sub) =>
              match norm (Expand.dist g), norm (Expand.dist f) with
              | .ok (g', subg), .ok (f', _) =>
                if equal g' f' then some ⟨F, "…", some ⟨Anti.integral f x, (steps.push check).push compare, F⟩, none⟩
                else some (refuse "integrate: the candidate … was rejected: its derivative simplifies to …")
      …
\end{minted}

Differentiate the candidate, normalize, expand both sides (Chapter~\ref{ch:parts} explains
why), and accept only on exact equality of the normal forms. The theorem is about what
acceptance means:

\begin{thmbox}[\lcn{cmdIntegrate_spec}]
  If \texttt{integrate(f, x)} returns $F$ without error, then normalizing $\frac{d}{dx} F$ gives
  some $g$, and $g$ and $f$ have the same normal form after expansion.
\end{thmbox}

\begin{minted}{lean}
theorem cmdIntegrate_spec (norm : Norm) {f : Expr} {x : String} {res : RuleResult}
    (h : (cmdIntegrate norm).apply (.fn "integrate" [f, .var x]) = some res) (hok : res.error = none) :
    ∃ g sub g' s₁ s₂, norm (D res.result x) = .ok (g, sub) ∧
      norm (Expand.dist g) = .ok (g', s₁) ∧ norm (Expand.dist f) = .ok (g', s₂)
\end{minted}

Nothing about the finder is assumed. The proof is a case analysis on the code: every branch that
returns without error passed the equality test. In \file{proofs/}, with the real semantics of
Chapter~\ref{ch:dderiv}, this becomes the statement a student would want:

\begin{minted}{lean}
theorem integrate_deriv (norm : Norm) … (ρ : EnvR)
    (hdiff  : ∀ g sub, norm (D res.result x) = .ok (g, sub) → evalD ρ (D res.result x) = evalD ρ g)
    (hleft  : …)   -- the expansion of g was sound at ρ
    (hright : …)   -- the expansion of f was sound at ρ
    : deriv (fx ρ x res.result) (ρ x) = evalD ρ f
\end{minted}

The hypotheses say that the three normalizations the checker ran were value-preserving at the
point $\rho$. Those normalizations are derivations shown in the notebook, each step with its
status; so the theorem's hypotheses are, exactly, ``the steps shown were sound at this point''.
That is why the command's status is \status{verified} while the finder's steps are
\status{checked}: the claim is the checker's, and it inherits the conditions of the derivative
rules it used.

\section{An inherited condition}

$\int x^a\,dx$ is accepted as $x^{a+1}/(a+1)$. Differentiating the candidate gives
$(a+1)x^a/(a+1)$, and the step that cancels $(a+1)/(a+1)$ is \rn{simp.collect-powers} of
Chapter~\ref{ch:exp}, conditional on a positive base --- here $a + 1 \neq 0$. So the check's
derivation carries the condition $a \neq -1$, the notebook shows it, and nobody had to write it
down. The finder's own note says ``assuming $n \neq -1$''; the checker's status \emph{proves} that
the assumption is where the answer's validity ends.

\begin{trybox}
\begin{minted}{text}
integrate(x^2 + sin(x), x)
integrate(exp(2*x), x)
integrate(x * cos(x^2), x)
integrate(x^a, x)
\end{minted}
  Open each derivation: the finder's steps are \status{checked}, then \rn{int.check} with the
  differentiation nested under it, then \rn{int.compare}. In the last, find the conditional
  step and read its condition.
\end{trybox}

\section{Exercises}

\begin{exercisebox}[Mean value theorem]
  State Rolle's theorem and derive the mean value theorem from it by subtracting the chord.
  Where in the proof is the hypothesis ``interval'' used?
\end{exercisebox}

\begin{exercisebox}[Seeing the substitution]
  Find $u$ for $\int \frac{\ln x}{x}\,dx$, $\int \sin x \cos x\,dx$ and
  $\int \frac{x}{\sqrt{1 - x^2}}\,dx$, and say which of the three the finder's \lc{trySubst}
  will find (the third needs a power with a non-integer exponent as $H'$).
\end{exercisebox}

\begin{exercisebox}[The constant]
  The engine's $\int x^{-1}\,dx$ is $\ln x$, valid for $x > 0$. The textbook's is $\ln|x|$.
  Explain, with the first theorem of this chapter, why the textbook's answer is two constants
  and not one.
\end{exercisebox}
